\scanpage{60}
\S~1 Coh. de De Rham ordinaire

\S~2 Tour d'horizon sur Coh. $l$-adiques. Motivations pour le désir d'une coh. $p$-adique en car. $p$.

\S~3 Tapis de Washnitzer-Monsky.

\S~4 Stratifications. Relations de W.M. avec connexions de Gauss-Manin sur coh. de De Rham

\S~5 \deleted{\unsure{Le}} Site cristallin et site stratifiant. Énoncé \deleted{de} d'un théorème et d'une conjecture en car. nulle.

\S~6 Calculs Čechistes. (car. quelconque).

\S~7 Comparaison de la Coh. de De Rham et cristalline (car. nulle).

\S~8 \deleted{\unsure{Considérations|Comparaison}} \inserted{Retour} sur le cas de la car. $p>0$

\noindent\rule{\linewidth}{.4pt}

\margin{\deleted{\uline{\unsure{Coh. topologique|Coh. stratifiée}}} \uline{Coh. stratifiante}}

6.1. Commençons à travailler sur le \deleted{\unsure{objet|site}}topos \uline{stratifiant} pour $X/S$. Il y a l'objet $\widetilde X$ de ce topos : il \uline{couvre l'objet final} (c'est ici qu'on utilise la restriction sur les objets du site stratifiant parmi ceux du site cristallin). Donc suite spectrale de Leray
\[
H^*(X_{\mathrm{Strat}},F)\Longleftarrow E^{pq}_2=H^p(\nu\rightsquigarrow H^q(\widetilde X^{\nu+1},F))
\]

Or on constate que
\[
\widetilde X^{\deleted{\unsure{p}}\nu+1}=\varinjlim_i\deleted{\unsure{\Delta_X}}\underbrace{\Delta^{(\nu)}_X(i)}_{\substack{\text{voisinage inf. $i$-ème}\\\text{de la diagonale dans}\\ X^{\nu+1}\unsure{\deleted{X'}|\longrightarrow X'}\text{, considéré comme \unsure{provoisin|provenant|provision|prosimplicial}}\\\text{de }X}}\deleted{\unsure{(X)}}
\]

D'autre part, \unsure{pour|\deleted{pour}} tt objet $(U,X')$ \deleted{\unsure{tel que}}
\diagram{\begin{tikzcd}[ampersand replacement=\&,row sep=tiny,column sep=tiny]
X \& U\arrow[l,hook']\arrow[dr]\\
\&\&X'
\end{tikzcd}}
du site stratifiant, on aura
% archive sheet number: 59

